% Figure 2 -- SoC integration architecture (v2 export, stage3_design-2026.09.29_ver2).
% GPIO_PADS is replaced by GPIO_PIN_SPLIT + eight ChipInventor Inout Pin blocks.
\documentclass[border=0pt]{standalone}
\usepackage{xcolor}
\usepackage{tikz}
\usetikzlibrary{arrows.meta,positioning,calc,fit,backgrounds}
% =====================================================================
%  Typography: IBM Plex superfamily (OFL, ships with TeX Live/MiKTeX
%  package "plex-otf"). Serif for reading, Sans for structure and
%  diagrams, Mono for code/logs, Plex Math for equations.
% =====================================================================
\usepackage{fontspec}
\usepackage{unicode-math}
\defaultfontfeatures{Ligatures=TeX}
\setmainfont{IBMPlexSerif}[
  Extension=.otf, Scale=0.92,
  UprightFont=*-Regular, BoldFont=*-SemiBold,
  ItalicFont=*-Italic, BoldItalicFont=*-SemiBoldItalic]
\setsansfont{IBMPlexSans}[
  Extension=.otf, Scale=0.95,
  UprightFont=*-Regular, BoldFont=*-SemiBold,
  ItalicFont=*-Italic, BoldItalicFont=*-SemiBoldItalic]
\setmonofont{IBMPlexMono}[
  Extension=.otf, Scale=0.95, HyphenChar=None,
  UprightFont=*-Regular, BoldFont=*-SemiBold,
  ItalicFont=*-Italic, BoldItalicFont=*-SemiBoldItalic]
\setmathfont{IBMPlexMath-Regular.otf}[Scale=0.95]
% Additional weights used for hierarchy
\newfontfamily\sanslight{IBMPlexSans}[Extension=.otf, Scale=0.95,
  UprightFont=*-Light, ItalicFont=*-LightItalic, BoldFont=*-Medium]
\newfontfamily\sansmedium{IBMPlexSans}[Extension=.otf, Scale=0.95,
  UprightFont=*-Medium, BoldFont=*-Bold]
\newfontfamily\sanscond{IBMPlexSansCondensed}[Extension=.otf, Scale=0.95,
  UprightFont=*-Regular, BoldFont=*-SemiBold,
  ItalicFont=*-Italic, BoldItalicFont=*-SemiBoldItalic]
\newfontfamily\sansdisplay{IBMPlexSans}[Extension=.otf,
  UprightFont=*-Bold, BoldFont=*-Bold]
\newfontfamily\sansthin{IBMPlexSans}[Extension=.otf,
  UprightFont=*-ExtraLight, BoldFont=*-Light]

% =====================================================================
%  RackSentinel visual system -- shared palette and TikZ vocabulary.
%  Loaded by main.tex (via rackstyle.sty) and by every standalone
%  figure source in figures/src/, so text and diagrams stay identical.
%
%  Colour carries meaning, never decoration:
%    navy    structure, headings, the processor
%    river   UART / data flow / links
%    teal    Normal state, PASS, verified
%    amber   Warm / Hot / overload (warning)
%    alarm   Critical / EPO only
%    mist    panels and quiet backgrounds
%    ink     body text
% =====================================================================
\definecolor{navy}{HTML}{15304B}
\definecolor{river}{HTML}{2A6FA8}
\definecolor{teal}{HTML}{1E8573}
\definecolor{amber}{HTML}{C9821A}
\definecolor{alarm}{HTML}{B5382C}
\definecolor{mist}{HTML}{EEF2F6}
\definecolor{mistdark}{HTML}{D9E1EA}
\definecolor{ink}{HTML}{24292F}
\definecolor{slate}{HTML}{5E6B78}
\definecolor{hair}{HTML}{B7C2CD}
% soft tints for state fills (still distinct in greyscale)
\definecolor{tealsoft}{HTML}{DCEFEA}
\definecolor{ambersoft}{HTML}{F7E7CC}
\definecolor{alarmsoft}{HTML}{F4D9D5}
\definecolor{riversoft}{HTML}{DDE9F4}

\newcommand{\RSTikzStyles}{%
\tikzset{
  >={Stealth[length=4.2pt,width=3.4pt]},
  every node/.style={font=\sffamily\small, text=ink},
  % blocks
  blk/.style={draw=navy, line width=0.7pt, fill=white, inner sep=3pt,
              align=center, minimum height=8mm},
  core/.style={blk, fill=navy, text=white, draw=navy},
  soft/.style={blk, fill=mist, draw=hair, line width=0.5pt},
  periph/.style={blk, fill=riversoft, draw=river},
  ext/.style={blk, draw=slate, dashed, line width=0.6pt, fill=white},
  % state chips
  sNormal/.style={fill=tealsoft, draw=teal},
  sWarn/.style={fill=ambersoft, draw=amber},
  sCrit/.style={fill=alarmsoft, draw=alarm},
  % wires
  uart/.style={draw=river, line width=1.1pt},
  gpio/.style={draw=navy, line width=0.9pt},
  bus/.style={draw=navy, line width=2.2pt},
  thin/.style={draw=slate, line width=0.5pt},
  lbl/.style={font=\sffamily\scriptsize, text=slate, inner sep=1.5pt},
  opt/.style={font=\ttfamily\scriptsize, text=river, inner sep=1.5pt},
  kicker/.style={font=\sffamily\tiny\bfseries, text=slate, inner sep=0pt},
}}

\RSTikzStyles
\begin{document}
\begin{tikzpicture}[x=1mm,y=1mm]
  % ---------------- clock/reset note ----------------
  \node[lbl, anchor=west, font=\sffamily\scriptsize\itshape] at (0,73) {\texttt{clk\_i}, \texttt{rst\_i} reach every block};

  % ---------------- core ----------------
  \fill[navy] (0,4) rectangle (38,68);
  \node[text=white, font=\sffamily\bfseries\normalsize, anchor=north] at (19,66) {RVBL-2 core};
  \node[text=white!80!navy, font=\sffamily\scriptsize, anchor=north] at (19,61.2) {Stage 2 multicycle, reused};
  \node[text=white, font=\sffamily\scriptsize, align=center, anchor=north, text width=34mm] at (19,54)
    {Control Unit (FSM)\\[1.3pt] PC, IR, Register File\\[1.3pt] ALU\\[1.3pt]
     \textbf{MULT} (Zmmul)\\[1.3pt] \textbf{CRC16} (Xicrc)\\[1.3pt] Imm.\ gen., LSU\\[1.3pt]
     Branch comparator\\[1.3pt] Datapath muxes};

  % ---------------- bus + decoder ----------------
  \draw[<->, navy, line width=1.8pt] (38,36) -- (50,36) node[midway, above=0.6mm, lbl] {bus};
  \draw[river, line width=1.1pt, fill=riversoft] (50,4) rectangle (72,68);
  \node[font=\sffamily\bfseries\small, text=navy, align=center] at (61,60) {Address\\decoder};
  \node[font=\sffamily\scriptsize, align=center, text width=19mm] at (61,38)
     {selects the block from the upper address bits;\\[2pt] steers \texttt{o\_Data} back to the core};
  \node[font=\sffamily\tiny\bfseries, text=river, align=center] at (61,9.5) {ONLY CORE BLOCK\\CHANGED IN STAGE 3};

  % ---------------- bus slaves ----------------
  \begin{scope}[every node/.append style={minimum width=40mm, minimum height=11.5mm, text width=38mm, align=center}]
    \node[blk] (imem) at (97.5,61) {\textbf{IMEM ROM}\enspace\texttt{0x0040\_0000}\\[-1pt]
          {\scriptsize \texttt{IMEM\_Data\_Center}, 275 words\\[-2pt] or \texttt{IMEM\_TESTFW}, 21 words}};
    \node[blk] (dmem) at (97.5,45) {\textbf{DMEM}\enspace\texttt{0x1001\_0000}\\[-1pt]{\scriptsize 8\,KB SRAM; 22 words of state}};
    \node[periph] (gpio) at (97.5,29) {\textbf{GPIO}\enspace\texttt{0xF000\_0000}\\[-1pt]{\scriptsize DATAOUT, DATAIN, DATADIR}};
    \node[periph] (uart) at (97.5,13) {\textbf{UART}\enspace\texttt{0xF100\_0000}\\[-1pt]{\scriptsize TXDATA, RXDATA, CONTROL\\[-2pt]\texttt{TX\_ON\_TXDATA\_WRITE = 0}}};
  \end{scope}
  \foreach \n in {imem,dmem,gpio,uart} \draw[<->, navy, line width=0.7pt] (72,0|-\n.west) -- (\n.west);
  \node[lbl, font=\sffamily\scriptsize\itshape, text=slate, align=left, anchor=west] at (119.6,63.5)
       {firmware lives only here:\\swap the ROM, keep\\the hardware};
  \draw[thin, -{Stealth[length=3pt]}] (119.2,61) -- (imem.east);

  % ---------------- GPIO_PIN_SPLIT + eight Inout Pins ----------------
  \node[periph, minimum width=20mm, minimum height=13mm, align=center] (split) at (138,29)
       {\textbf{GPIO\_PIN\_}\\[-2pt]\textbf{SPLIT}\\[-1pt]{\scriptsize bus ↔ 8 pins}};
  \draw[gpio,->] (gpio.east|-0,31.5) -- (split.west|-0,31.5) node[midway, above=-0.2mm, lbl, font=\sffamily\tiny] {out, oe};
  \draw[gpio,<-] (gpio.east|-0,26.5) -- (split.west|-0,26.5) node[midway, below=-0.2mm, lbl, font=\sffamily\tiny] {in};
  \foreach \k in {0,...,7} {
    \pgfmathsetmacro{\px}{127.8+\k*2.75}
    \draw[navy, line width=0.45pt] (\px+1,35.5) -- (\px+1,40);
    \draw[navy, line width=0.55pt, fill=white] (\px,40) rectangle ++(2,4.4);
  }
  \node[lbl, anchor=south, align=center, text=navy, font=\sffamily\tiny] at (138.8,44.5)
       {8 Inout Pin blocks, \texttt{D\,?\,C\,:\,1'bz}};
  \node[font=\ttfamily\scriptsize\bfseries, text=navy, anchor=south] at (138.8,47.6) {pins\_io\_0 … pins\_io\_7};

  % ---------------- UART pins ----------------
  \draw[uart,->] (uart.east|-0,15.5) -- (154,15.5)
        node[pos=0.55, above=0.2mm, font=\ttfamily\scriptsize\bfseries, text=river] {tx\_o};
  \draw[uart,<-] (uart.east|-0,10.5) -- (154,10.5)
        node[pos=0.55, below=0.2mm, font=\ttfamily\scriptsize\bfseries, text=river] {rx\_i};

  % ---------------- external stimulus ----------------
  \node[ext, minimum width=22mm, minimum height=62mm, anchor=north east] (emu) at (178,68) {};
  \node[font=\ttfamily\small\bfseries, text=slate] at (167,62.5) {./emu};
  \node[font=\sffamily\scriptsize, text=slate, align=center, text width=20mm] at (167,57) {test host,\\outside the chip};
  \node[font=\sffamily\scriptsize, text=slate, align=center, text width=20mm] at (167,36)
       {\texttt{-d 0xF0}\\[2pt] drives P0\\with \texttt{-i}\\[2pt] checks P7–P4\\with \texttt{-o}};
  \node[font=\sffamily\scriptsize, text=slate, align=center, text width=20mm] at (167,13)
       {sends \texttt{-t4}, \texttt{-t2}\\expects \texttt{-r}};
  \draw[gpio,<->, draw=slate, dashed] (149.8,42.2) -- (emu.west|-0,42.2);
  \draw[uart,dashed] (154,15.5) -- (emu.west|-0,15.5);
  \draw[uart,dashed] (154,10.5) -- (emu.west|-0,10.5);

  % ---------------- legend + note ----------------
  \begin{scope}[shift={(0,-5.4)}]
    \fill[riversoft] (0,0) rectangle (4,2.6); \draw[river, line width=0.6pt] (0,0) rectangle (4,2.6);
    \node[lbl, anchor=west] at (4.8,1.3) {new or changed in Stage 3};
    \draw[navy, line width=0.7pt, fill=white] (40,0) rectangle (44,2.6);
    \node[lbl, anchor=west] at (44.8,1.3) {existing block};
    \draw[slate, dashed, line width=0.6pt] (66,0) rectangle (70,2.6);
    \node[lbl, anchor=west] at (70.8,1.3) {external test host};
    \node[lbl, anchor=east, font=\sffamily\scriptsize\itshape, text=navy] at (178,1.3)
       {Only the decoder knows the memory map; GPIO and UART decode bits 11:2.};
  \end{scope}
\end{tikzpicture}
\end{document}
